\documentclass{aediroum}

\newenvironment{instance}{\flushright\itshape\scriptsize}

\title{Cahier de positions}
\date{17 septembre 2025}

\begin{document}
\maketitle

Le présent document regroupe l'ensemble des positions de l'AÉDIROUM. Il est utilisé par toute personne faisant de la représentation au nom de l'association étudiante. Ce document doit être mis à jour régulièrement afin de bien refléter les positions des membres de l'AÉDIROUM.

Ce document a été créé à la session d'hiver 2011 par l'exécutif en place afin de recenser les opinions de ses membres, abrogé en 2016, puis reconstitué en 2021. Les positions sont organisées en deux sections selon leur champ d'application.

\section{Affaires externes}\label{sec:affaires-externes}

\begin{enumerate}
	\item \marginnote{Frais de scolarité} Que l'AÉDIROUM travaille afin que le système d'éducation supérieure québécois:
		\begin{itemize}
			\item Soit également accessible indépendamment de la situation financière des étudiants et de leurs proches;
			\item Permette à un maximum d'étudiants compétents de compléter leurs projets d'études et d'emploi.
		\end{itemize}
		\begin{instance}
			Adopté: [CA-2022-03-18]
		\end{instance}

	\item Que l'AÉDIROUM s'oppose à toute hausse des frais de scolarité, en perspective de la gratuité scolaire.
		\begin{instance}
			Adopté: [CA-2022-03-18]
		\end{instance}

	\item \marginnote{Environnement} Que l'AÉDIROUM mette de l'avant et défende les initiatives environnementales.
		\begin{instance}
			Adopté: [CA-2022-03-18]
		\end{instance}

	\item \marginnote{Transport en commun} Que l'AÉDIROUM appuie des démarches en vue d'avoir un tarif réduit pour tous les étudiants à temps plein.
		\begin{instance}
			Adopté: [CA-2022-03-18]
		\end{instance}

	\item \marginnote{Association de campus} Que la FAÉCUM facilite l'accès à ses documents, notamment en permettant de rendre public ses documents d'intérêt général, dont ses règlements généraux et sa politique d'accès à l'information.
		\begin{instance}
			Adopté: [CA-2022-03-18]
		\end{instance}

        \item Que l'AÉDIROUM demande à la FAÉCUM de permettre à tous ses membres, individuels et associatifs, de consulter et d'obtenir une copie de son budget complet et de ses états financiers en tout temps et sans délai, et ce dans un format humainement compréhensible, le tout dans un objectif de transparence.
		\begin{instance}
			Adopté: [AG-2025-09-17]
		\end{instance}

        \item Que l'AÉDIROUM demande à la FAÉCUM de rendre son processus électoral plus ouvert et transparent, en encourageant activement la candidature de personnes issues d'une pluralité d'associations au sein de son bureau exécutif et de son conseil d'administration.
		\begin{instance}
			Adopté: [AG-2025-09-17]
		\end{instance}

        \item Que l'AÉDIROUM demande à la FAÉCUM de permettre à ses personnes membres de s'en désinscrire à titre individuel et de cesser d'en payer les cotisations, à condition qu'elles ne soient pas membres d'une association étudiante affiliée à la Fédération.
		\begin{instance}
			Adopté: [AG-2025-09-17]
		\end{instance}

        \item Que l'AÉDIROUM appuie les membres de la FAÉCUM dans une démarche qui vise à donner plus de voix à ses membres à titre individuel, y compris au travers de regroupements étudiants, ainsi qu'à renforcer le caractère sain et démocratique des instances de la Fédération et la souveraineté de ses membres dans celles-ci.
		\begin{instance}
			Adopté: [AG-2025-09-17]
		\end{instance}

	\item \marginnote{Discrimination} Que l'AÉDIROUM se positionne contre le sexisme, l'homophobie, la transphobie, l'hétéronormativité et toutes formes de discrimination basées sur le sexe, le genre ou l'orientation sexuelle. Que l'AÉDIROUM prenne acte des inégalités persistantes entre les hommes et les femmes et qu'elle s'oppose à leurs perpétuations.
		\begin{instance}
			Adopté: [CA-2022-03-18]
		\end{instance}

	\item \marginnote{Démocracie} Que l'AÉDIROUM appuie le principe de démocratisation des institutions d'enseignement dans une perspective d'autogestion. Que l'AÉDIROUM se positionne pour l'élection du conseil d'administration et du rectorat de l'Université de Montréal par la communauté universitaire.
		\begin{instance}
			Adopté: [CA-2022-03-18]
		\end{instance}

	\item \marginnote{Aide financière} Que l'AÉDIROUM soit pour un régime d'aide financière adéquat ayant pour but d'éliminer l'endettement étudiant et d'assurer la satisfaction des besoins fondamentaux.
		\begin{instance}
			Adopté: [CA-2022-03-18]
		\end{instance}

	\item \marginnote{Entreprises privées} Que l'AÉDIROUM soit pour un réseau d'éducation public libre de toute ingérence de l'entreprise privée au niveau de la gouvernance des institutions publiques, et que l'AÉDIROUM condamne la culture de sous-traitance abusive dans les institutions publiques d'éducation.
		\begin{instance}
			Adopté: [CA-2022-03-18]
		\end{instance}

        \item Que l'AÉDIROUM s'oppose à toute délocalisation des services informatiques de l'Université à des compagnies externes et qu'elle encourage l'Université à tirer parti de ses compétences à l'interne pour le développement d'outils numériques libres, indépendants et ouverts.
                \begin{instance}
                        Adopté: [AG-2025-09-17]
                \end{instance}

	\item Que l'AÉDIROUM soit contre toute forme de mondialisation qui entérine la prédominance du profit sur le bien-être de la population.
		\begin{instance}
			Adopté: [CA-2022-03-18]
		\end{instance}

	\item \marginnote{Inclusivité} Que l'AÉDIROUM dénonce toute forme de discrimination ciblant l'identité ou l'expression de genre.
		\begin{instance}
			Adopté: [CA-2022-03-18]
		\end{instance}

	\item Que l'AÉDIROUM appuie le regroupement étudiant de l'Université de Montréal \emph{l'Alternative} dans ses démarches pour une université inclusive.
		\begin{instance}
			Adopté: [CA-2022-03-18]
                        Amendé: [AG-2025-09-17]
		\end{instance}

	\item Que l'AÉDIROUM invite l'Université de Montréal~:
		\begin{enumerate}
			\item à modifier sa Politique contre le harcèlement afin d'interdire explicitement la discrimination envers l'identité ou l'expression de genre, et à promouvoir l'acceptation des personnes trans*~;
			\item à faciliter les démarches permettant de concilier l'identité de genre des étudiants et leur statut administratif, notamment en permettant aux étudiants d'utiliser leur nom, leur formule d'appel et leur mention de sexe préférés sur leur carte étudiante, sur leur horaire, sur les listes de classe, sur l'environnement StudiUM, dans leur courriel institutionnel, et dans tout autre contexte où il est possible de le faire~;
			\item à créer des espaces plus accessibles pour les personnes trans*, notamment en incluant des toilettes neutres sur ses plans de construction ou de rénovation futurs ou en affichant les toilettes à une place comme neutres~;
			\item à faciliter, dans la mesure du possible, l'accès aux soins de santé que requièrent les personnes trans* à la clinique universitaire, et à éduquer le personnel soignant pour qu'il sache bien gérer les personnes trans*~;
			\item à diffuser, tant aux étudiant-e-s qu'aux employé-e-s, des informations sur l'identité et l'expression de genre et sur les politiques de l'Université visant les personnes trans*.
		\end{enumerate}
		\begin{instance}
			Adopté: [CA-2022-03-18]
		\end{instance}

	\item \marginnote{Mode de scrutin} Que l'AÉDIROUM soutienne la cause du SENSÉ dans ses démarches pour un nouveau mode de scrutin à mi-chemin entre la représentation circonscriptions et le vote proportionnel mixte. Ce mode permet aux partis politiques émergents d'être moins pénalisés par le système de circonscriptions et surtout de donner un pourcentage de pouvoir à chaque parti qui correspond au pourcentage réel de la population derrière chaque parti.
		\begin{instance}
			Adopté: [CA-2022-03-18]
		\end{instance}
\end{enumerate}

\subsection{Conflits au Proche-Orient}

\begin{enumerate}
	\item \marginnote{Apartheid} Que l'AÉDIROUM reconnaisse l'apartheid imposé par l'État d'Israël.
		\begin{instance}
			Adopté: [AG-2024-11-13]
		\end{instance}
	
	\item \marginnote{Génocide} Que l'AÉDIROUM reconnaisse le génocide perpétré par l'État d'Israël contre le peuple palestinien.
		\begin{instance}
			Adopté: [AG-2024-11-13]
		\end{instance}
	
	\item Que l'AÉDIROUM soutienne le peuple palestinien dans la défense de son territoire et de sa culture.
		\begin{instance}
			Adopté: [AG-2024-11-13]
		\end{instance}
	
	\item \marginnote{Liens académiques} Que l'AÉDIROUM demande à l'Université de Montréal d'interrompre ses liens académiques avec toute université israëlienne.
		\begin{instance}
			Adopté: [AG-2024-11-13]
		\end{instance}
	
	\item \marginnote{Boycott} Que l'AÉDIROUM encourage ses membres ainsi que l'Université de Montréal à boycotter l'État d'Israël.
		\begin{instance}
			Adopté: [AG-2024-11-13]
		\end{instance}
	
	\item \marginnote{OTAN} Que l'AÉDIROUM dénonce le soutien de l'OTAN et de ses membres au génocide perpetré par l'État d'Israël contre le peuple palestinien.
		\begin{instance}
			Adopté: [AG-2024-11-13]
		\end{instance}
	
	\item \marginnote{Cessez-le-feu} Que l'AÉDIROUM demande un cessez-le-feu immédiat et le retrait des troupes armées israéliennes en Palestine.
		\begin{instance}
			Adopté: [AG-2024-11-13]
		\end{instance}

	\item Que l'AÉDIROUM demande un cessez-le-feu immédiat et le retrait des troupes armées israéliennes au Liban.
		\begin{instance}
			Adopté: [AG-2024-11-13]
		\end{instance}
\end{enumerate}

\section{Affaires internes}\label{sec:affaires-internes}

\begin{enumerate}
	\item \marginnote{Organisation des cours} Que l'AÉDIROUM s'assure de la tenue d'une rétroaction de mi-session toutes les sessions d'automne et d'hiver.
		\begin{instance}
			Adopté: [CA-2022-03-18]
		\end{instance}

        \item Que l'AÉDIROUM s'oppose à l'assignation de devoirs ou travaux à remettre pendant les périodes d'examens intra ou finaux.
                \begin{instance}
                        Adopté: [AG-2024-02-07]
                \end{instance}

	\item \marginnote{Logiciels libres} Que l'AÉDIROUM préconise l'utilisation de logiciels libres et de formats ouverts.
		\begin{instance}
			Adopté: [CA-2022-03-18]
		\end{instance}

        \item Que l'AÉDIROUM encourage l'utilisation des logiciels libres sur le plan individuel et qu'elle demande l'appui de l'université dans ses démarches
		\begin{instance}
			Adopté: [AG-2025-09-17]
		\end{instance}

	\item \marginnote{Politique linguistique} Que l'AÉDIROUM s'assure que tous les cours et le matériel d'apprentissage soient disponibles en français au premier cycle.
		\begin{instance}
			Adopté: [CA-2022-03-18]
		\end{instance}

        \item \marginnote{Budget} Que l'AÉDIROUM permette à ses membres de consulter et d'obtenir une copie de son budget complet et de ses états financiers en tout temps et sans délai, et ce dans un format humainement compréhensible, le tout dans un objectif de transparence.
		\begin{instance}
			Adopté: [AG-2025-09-17]
		\end{instance}

	\item Que l'AÉDIROUM défende le choix des étudiants en ce qui a trait au choix de la langue d'écriture de leur thèse de doctorat ou de leur mémoire de maîtrise.
		\begin{instance}
			Adopté: [CA-2022-03-18]
		\end{instance}

	\item \marginnote{Identité} Que l'AÉDIROUM reconnaisse l'identité de genre déclarée de tout-e-s ses membres dans le cadre de toutes ses activités, peu importe leur statut légal, administratif, chirurgical ou autre.
		\begin{instance}
			Adopté: [CA-2022-03-18]
		\end{instance}

	\item Que l'AÉDIROUM se positionne au niveau départemental pour qu'on puisse utiliser son nom usuel sur les travaux et les examens.
		\begin{instance}
			Adopté: [CA-2022-03-18]
		\end{instance}

        \item \marginnote{Autres} Que l'AÉDIROUM remercie Robin Milosz.
                \begin{instance}
                        Adopté: [AG-2024-02-07]
                \end{instance}
\end{enumerate}

\end{document}
